\subsection{Stochastic vs Deterministic Channel Modeling}
Channel modeling can be broadly classified into two categories, namely stochastic and deterministic channel modeling.
Stochastic models rely on satistcial distributions to represent large-scale effects such as path loss and shadowing as well as small-scale fading through multipath clusters \cite{tr38901}.
These models are computationally efficient and widely adopted in system-level simulators to evaluate average network performance across standardized deployment scenarios \cite{tr38901,5glena}.

Eventhough they are suitable for lage-scale evalution, they lack site-specific realistic geometric information about the environment \cite{alsamman_mmwave,han_thz}.
For instance, a stochastic model cannot distinguish between a line-of-sight path along an open space and blocked path behind a building as both are represented through probalistic parameters.
In constrast, deterministic modeling computes the actual electromagnetic interactions within a given 3D environment.

While full-wave solutions based on Maxwell's equations are computationally infeasible, high-frequency approximations such as ray-tracing provide a pratical alternative \cite{alsamman_mmwave,han_thz}.
Ray-tracing bridges the gap between statistical averages and physical reality by resolving propagation paths including reflections, diffrations and blockages.
This is especially important for 6G research as they use mmWave and THz frequency bands, where propagation is highly directional and sensitive to enviorment geometry \cite{uwaechia_mmwave,han_thz}.
